\section{Basic Concepts of Thermodynamics}

\subsection{Systems and Properties}
\begin{itemize}
    \item \textbf{Open System}: Exchanges both matter and energy with surroundings.
    \item \textbf{Closed System}: Exchanges energy but not matter.
    \item \textbf{Isolated System}: Exchanges neither matter nor energy.
    \item \textbf{Intensive Properties}: Independent of the amount of substance (e.g., $T, P, \rho$, molar volume).
    \item \textbf{Extensive Properties}: Dependent on the amount of substance (e.g., $m, V, U, H, S, G$).
\end{itemize}

\subsection{State Functions and Processes}
\begin{itemize}
    \item \textbf{State Function}: A property that depends only on the current state, not the path (e.g., $P, V, T, U, H, S, G$).
    \item \textbf{Thermodynamic Process}: The manner in which a system changes from an initial state to a final state.
\end{itemize}

\subsection{Gases}
\textbf{Ideal Gas Law:}
\begin{equation}
    PV = nRT \quad \text{or} \quad PV_m = RT
\end{equation}
\textbf{Real Gas (Van der Waals Equation):}
\begin{equation}
    \left( P + \frac{a}{V_m^2} \right)(V_m - b) = RT
\end{equation}
Where:
\begin{itemize}
    \item $a/V_m^2$: correction for intermolecular attractive forces.
    \item $b$: volume excluded by intermolecular repulsive forces.
\end{itemize}

\textbf{Dalton's Law of Partial Pressures:}
\begin{equation}
    P = \sum P_i, \quad P_i = x_i P, \quad x_i = \frac{n_i}{n}
\end{equation}

\subsection{Zeroth Law of Thermodynamics}
If two systems (A and B) are each in thermal equilibrium with a third system (C), then they are in thermal equilibrium with each other ($T_A = T_B = T_C$).
